\section{An explicit computable size ceiling}\label{sec:quantitative}

We give an ordinary computable upper bound for the cardinality of one
nonrepresentable lattice in the affine-chart family.  The bounded-rank
exclusion in \Cref{sec:bounded} used a sufficiently-general-subgroup
theorem with an unspecified uniform threshold.  To choose explicit finite
parameters here, we first prove a quantitative substitute: a finite subgroup
avoiding proper closed sets of a specified computable format lies between
the derived group and the full fixed group of a Steinberg map.  We then
bound the pruning and protected-path parameters and count the resulting
lattice.

These bounds concern the negative affine family, independently of the
weighted-template counting bound and of any hypothetical representation
decider.  They provide the fixed initial segment of lattice tables used in
\Cref{sec:minimum}.

\subsection{A uniform expanding function}

\begin{definition}[Affine ring format]\label{q:format}
For an integer $x\ge10^9$, an affine ring formula has \emph{format at most
$x$} if it has at most $x$ occurrences in total of polynomial-equality atoms
and the operations $\neg,\land,\lor,\exists,\forall$, at most $x$ variables,
and polynomial atoms of total degree at most $x$.  Occurrences are counted in
the formula tree, with a separate variable for each quantifier binder.  The
variable count includes the entire ambient affine coordinate list, even
unused coordinates.  Coefficients are arbitrary field elements, and a
polynomial is counted through its atom and degree, rather than its expanded
arithmetic expression.

Let $a(x)$ be the least positive simultaneous bound, over all algebraically
closed fields and formulas of format at most $x$, for both the format of an
equivalent quantifier-free formula in the same ambient coordinates and the
number and degrees of polynomial equations defining its Zariski closure.
\end{definition}

\begin{lemma}\label{q:acf-computable}
The integer $a(x)$ exists, is nondecreasing, and is total computable.  This
holds uniformly over the characteristic and the specialization of every
coefficient.
\end{lemma}

\begin{proof}
For fixed $x$, canonically rename the coordinate and bound variables.  There
are finitely many possible formula trees.  Replace each polynomial atom by
the polynomial containing every allowed monomial of degree at most $x$, with
an independent coefficient parameter.  This produces finitely many schemes.
Their parameter lists are finite and computable, although the specialized
coefficients need not have effective descriptions.  These parameters belong
to the tests of uniformity, not to the variable count of the specialized
input formula.

Effective quantifier elimination in the theory ACF, without fixing the
characteristic, supplies a quantifier-free formula for each scheme
\cite{CohenMahboubi2012}.  Write it as a finite disjunction of conjunctions.
For a conjunction $f_i(z)=0$, $g_j(z)\ne0$, introduce one inverse variable $u$
and the ideal
\[
 I=\bigl(f_i(z),\ u\textstyle\prod_jg_j(z)-1\bigr).
\]
The Nullstellensatz gives the closure of that conjunction as $V(I\cap k[z])$ \citep[Theorem~10.34.1]{StacksNullstellensatz}.  An elimination
Gr\"obner basis gives equations for this closure.  The closure of the original
finite union is the union of these closures, which again has a finite
equation list, for example from the product of their defining ideals.
Characteristic-independent Gr\"obner bounds are available
\cite{Dube1990}; we give also the uniformity argument needed to compute the
\emph{least} simultaneous bound.

For given bounds on variables, generators, and degrees, the assertion that a
tuple of generators has a Gr\"obner basis with certificates of size at most
$b$ is first order in its coefficients.  The certificates consist of bounded
lists of polynomials, bounded ideal combinations in both directions, choices
of leading monomials from finite lists, and Buchberger reductions or standard
representations.  The required monomial comparisons for a fixed elimination
order form finite tests; polynomial identities become coefficient equations.
We can make these assertions increasing in $b$.  Every individual tuple in
every algebraically closed field has some certificate.  If no uniform $b$
existed, the Compactness Theorem \citep[p.~26]{Marker2020}, applied to ACF with constants for the coefficients and
the negations of all these increasing assertions, would give a tuple with no
certificate at all.  Thus a uniform $b$ exists.  First-order ACF tests of
successive $b$ find one.

In such a certified elimination basis, the members involving only the
retained coordinates define the projected closure.  Independence of the
eliminated variables is itself a finite coefficient test.  Existentially
quantifying a bounded certified basis therefore gives an effectively found
formula $C_\phi(c,z)$ for the closure of the fiber defined by $\phi(c,z)$.
The coefficient specialization precedes the certificate test.  In
particular, this constructs each specialized fiber's closure, rather than
specializing the closure of the total parameter family.  Vanishing leading
coefficients, empty pieces, and empty equation lists are handled by finite
branches.

For a candidate bound $A$, enumerate the finitely many quantifier-free
templates $Q(d,z)$ of format at most $A$ in the same coordinates.  For each
input scheme, the assertion
\[
 \forall c\ \bigvee_Q\exists d\ \forall z\,
       \bigl(\phi(c,z)\longleftrightarrow Q(d,z)\bigr)
\]
tests the quantifier-free part of the bound.  Let $P_1(e,z),\ldots,P_A(e,z)$
be dense polynomials of degree at most $A$ with coefficient tuple $e$.
The additional assertion
\[
 \forall c\ \exists e\ \forall z\,
 \left(C_\phi(c,z)\longleftrightarrow
                       \bigwedge_{i=1}^{A}P_i(e,z)=0\right)
\]
tests the closure bound; zero polynomials pad shorter lists.  Conjoin these
tests over the finitely many schemes and ambient dimensions.  Validity of
the resulting sentence in all algebraically closed fields is decidable:
after effective elimination a closed sentence is a Boolean combination of
integer equalities, whose behavior is checked in characteristic zero and at
the finitely many prime divisors of the nonzero integers occurring.  Testing
$A=1,2,\ldots$ consequently finds exactly $a(x)$.  The nested input classes
make $a$ nondecreasing.  Finally, B\'ezout bounds give elementary bounds on
the number of components and their sum of degrees from these equation
bounds.  All these statements allow ambient affine matrix coordinates,
including equations specifying the group itself.
\end{proof}

\begin{definition}[The function $J$]\label{q:J}
Let $R(c,h)$ be the least positive set size forcing a monochromatic
$h$-element subset in every $c$-coloring of its unordered pairs of distinct
elements.  Let $\operatorname{HJ}(a,c,s)$ be the least positive word length
forcing a monochromatic $s$-slot word family over an $a$-letter alphabet in
every $c$-coloring.  Here some positions have fixed letters, possibly no
positions do, and the others are partitioned into $s$ disjoint nonempty
sets; each set carries one independently arbitrary common letter.  Put
\begin{align}
 J(x)=1+\max\bigl(&\{2^{2^{x!}},a(x)\}
     \cup\{R(c,h):1\le c,h\le x\}\notag\\
     &\cup\{\operatorname{HJ}(a,c,s):1\le a,c,s\le x\}\bigr).
 \label{q:J-formula}
\end{align}
The notation $J^d$ denotes iteration, not an ordinary power.
\end{definition}

The finite Ramsey and multidimensional Hales--Jewett theorems
\cite{Ramsey1930,HalesJewett1963} ensure that these least bounds exist.
For each parameter tuple they can be found, in principle, by finite coloring
tests at successively larger sizes.  Thus $J$ is nondecreasing and total
computable.  Its double-exponential factorial term ensures that an additional
iteration absorbs the elementary polynomial and exponential costs used below.
No evaluation of these finite searches is required.

\subsection{Formats for bounded root systems}

We make the genericity bound quantitative in the affine coordinates of the
adjoint action.  The conventions for positive standard point maps, including
the parameter $P$ for the special twisted types, are those of
\Cref{sec:bounded}.

\begin{lemma}[Root words and fixed points]\label{q:root-formats}
Let $10^9\le x\le y$ be integers, and let $G$ be split pinned, adjoint
simple of rank at most $x$ over an algebraically closed field of positive
characteristic.  Its adjoint representation embeds $G$ as a closed subgroup
of $\operatorname{GL}_d$, where $d=\dim G=O((x+10)^2)$.
For rational simple modules of the simply connected cover, of dimension
and highest-weight coefficients at most $y$, graphs of the induced
projective actions have format at most $J(y)$.  Stabilizers of a subalgebra
or of at most $y$ subspaces of the
corresponding matrix algebra, and normalizers of projective subgroups of
order at most $y$, have format at most $J^3(y)$.

For a positive standard map of parameter $P$, the fixed-point test has a
Boolean description with at most $J^5(y)$ atoms and degrees at most
$J^5(y)P$.  A proper closed subset of $G$ of format at most $y$ therefore
contains at most
\begin{equation}\label{q:proper-fixed-count}
                 J^{10}(y)P^{\dim G-1}
\end{equation}
fixed points.  For $P>J(y)$ the derived simple fixed group has order at least
$P^{\dim G}/(10y)^4$.  The same bounds apply to the other adjoint types in
the stated rank range.
\end{lemma}

\begin{proof}
\emph{The embedding.}
Every root is primitive in the character lattice of the adjoint torus,
which is the root lattice: it is Weyl-conjugate to a simple root.  Its
differential is therefore nonzero in every characteristic.  An element
central in the Lie algebra has no root-space components, by commuting with
toral directions.  Its remaining toral component vanishes because the
simple-root differentials form a basis.  Thus the adjoint kernel has zero
Lie algebra.  It is finite reduced, is central by connectedness, and is
trivial because $G$ is adjoint.  The separable isomorphism theorem for
algebraic groups identifies $G$ with its closed adjoint image.  The image
has determinant one.  We use a Frobenius-standard basis.

\emph{Word and stabilizer formats.}
Bruhat decomposition and the rank-one torus formulas cover both simply
connected and adjoint points by root-element words of length less than
$1000(y+10)^4$.  This includes zero parameters and words for Weyl
representatives.  Root matrices on the modules in question have polynomial
entries of bounded degree: a coefficient of $t^h$ shifts weight by $h\alpha$,
and the weights lie in the Weyl polytope.  Coroot heights in the irreducible
root lists are less than $100(y+10)$, so pairing with the coroot bounds the
rootwise weight range.  The number of word choices needed is at most
\begin{equation}\label{q:word-choices}
             \bigl(1000(y+10)^2\bigr)^{1000(y+10)^4+1}.
\end{equation}
Use entries of intermediate products as variables and the coefficients of
the root maps as parameters.  Describing a group element and its action by
the same word gives the asserted $J(y)$ format, using invertible matrices
modulo scalar matrices for a projective action.  The same bound also
dominates Weyl dimensions for the indicated highest-weight coefficients, by
the Weyl dimension formula.  Conditions that a matrix preserve a subspace,
a subalgebra, a collection of at most $y$ subspaces, or normalize a
projective group with at most $y$ listed elements add only bounded lists of
variables and equations.  Allowing matrices and translates as parameters
gives the $J^3(y)$ bound.  Images, preimages, and intersections combine these
formulas; closures use \Cref{q:acf-computable}.

\emph{Point maps and the fixed-point estimate.}
Pinned diagram maps and special isogenies in their fixed characteristics
have bounded graph format by the same words: their root parameters are
sent to $t$ or $t^p$.  Split field powers, including their inverses on
points, only twist the coefficients of sets.  Inverses of special maps use
also an inverse field power.  Images of bounded sets under these point
maps thus cost only a few further iterations of $J$.

For a positive standard map, first eliminate the bounded combination of
inner, pinned diagram, or special maps, whose graph has format at most
$J^2(y)$, and only then substitute the single split field power.  This
gives at most $J^5(y)$ atoms of degree at most $J^5(y)P$ in the fixed-point
test.  On a conjunction piece, the inequalities specify an open subset.
Since the fixed set is finite, only isolated components after the
polynomial cuts can contribute.  On a variety of dimension $d$ and sum of
component degrees $\delta$, cuts of degree at most $E$ contribute at most
$\delta E^d$ such points: choose at most $d$ generic linear combinations
of the cuts retaining the finitely many points as isolated points, and
apply affine B\'ezout for locally closed sets \citep[Definition~3, Proposition~2.9 and Theorem~2.14]{PardoSebastian2020}.  The generic-combinations step supplies the exponent $d$ locally.  For a proper subset of $G$, $d\le\dim G-1$.
The degree bound obtained from $a(y)$, the number of Boolean pieces, and
the factor $(J^5(y))^{\dim G}$ are absorbed by $J^{10}(y)$, proving
\eqref{q:proper-fixed-count}.  The ordinary order and center formulas give
the stated lower bound for the derived simple when $P>J(y)$.
\end{proof}

\begin{proposition}[Quantitative genericity]\label{q:genericity}
Let $x\ge10^9$, and let $G$ be adjoint simple of rank at most $x$ over an
algebraically closed field of characteristic $p>0$.  If a finite subgroup
$\Gamma\le G$ is not contained in any proper closed subset of affine ring
format at most $J^{1000}(x)$, then there is a positive-height standard point
map $\Psi$ such that
\begin{equation}\label{q:sandwich}
                   (G^\Psi)'\le\Gamma\le G^\Psi.
\end{equation}
Proper closed subsets in this hypothesis need not be irreducible.
\end{proposition}

\begin{proof}
Choose a fundamental simple module of the simply connected cover of $G$.
Its dimension $n$ is at most $J(x)$.  Its weights lie in one coset of the
root lattice, so the center acts by a scalar character, giving an
irreducible projectively point-faithful action of $G$.

\emph{Claim 1. The matrix lifts of every nontrivial normal subgroup of
$\Gamma$ span the full matrix algebra, and $\Gamma$ is almost simple.}
This connected action preserves no nontrivial tensor splitting individually.
Indeed the projective actions on the factors are algebraic, as can be seen
from their conjugation actions on the factor matrix algebras.  Pull back
$\operatorname{SL}\to\operatorname{PGL}$ along each map from the simply
connected source.  The reduced identity group of the pullback still
surjects onto that source: the full map is finite surjective and has only
finitely many components.  Its kernel is central of multiplicative type.
Its connected solvable radical has trivial image and hence dimension zero,
so it is semisimple.  The resulting central isogeny onto the simply
connected source is an isomorphism, providing a linear lift.  Perfection
leaves no scalar character in the tensor comparison.  Irreducibility would
make both tensor factors nontrivial irreducible modules, but their dominant
highest weights cannot add to a single fundamental weight.

An invariant unital matrix subalgebra is accordingly scalar or full.  Its
radical would otherwise give a proper invariant annihilator.  If it is
semisimple, connectedness fixes its central idempotents, so irreducibility
leaves one matrix factor.  A proper nonscalar such factor would yield the
excluded tensor splitting.

We may use $y=J^2(x)$ in the stabilizer tests of
\Cref{q:root-formats}; it covers $n^2$ and all the lists of matrix data
needed here.  Their stabilizers and the bounded unions of translates used
below have format well within $J^{1000}(x)$.  Avoidance implies
$\Gamma\ne1$.  For every nontrivial normal subgroup of $\Gamma$, the
algebra spanned by its matrix lifts has a bounded-format stabilizer
containing $\Gamma$, so is $G$-invariant.  It is nonscalar by projective
faithfulness, and hence is the full matrix algebra.

If such a normal subgroup were abelian, its distinct projective elements
would have lifts with distinct scalar conjugation characters.  Equal
characters would make the quotient of two lifts centralize the full span,
forcing equal projective elements.  Independence of the corresponding
eigenvectors bounds the order by $n^2$.  Its normalizer is then a bounded
proper closed subgroup: a finite nontrivial subgroup normal in connected
adjoint simple $G$ would be central, which is impossible.  This contradicts
avoidance.

A minimal normal subgroup is therefore a power of a nonabelian simple
group.  Lifts of its factors commute linearly across factors, by
perfection, and their algebras jointly generate the full matrix algebra.
A radical in one algebra would generate a nilpotent ideal in that full
algebra, so vanishes; the center is scalar for the same commuting-generation
reason.  Each factor algebra is consequently a full matrix algebra, and
their tensor product is simple and maps onto the full matrix span.  These
nontrivial tensor factors are permuted by $\Gamma$.  The bounded stabilizer
of their collection must contain $G$, whose connectedness fixes each
factor.  Tensor indecomposability leaves exactly one factor.  Denote the
resulting simple normal subgroup by $S$.  Fullness and perfection give it
trivial projective centralizer, and conjugation identifies
\begin{equation}\label{q:almost-simple}
                       S\le\Gamma\le\operatorname{Aut}(S).
\end{equation}

\emph{Claim 2. The simple normal subgroup $S$ is of Lie type in
characteristic $p$, of rank at most $100(n+1)$ and parameter
$P>J^{100}(x)$.}
The classification of finite simple groups, together with the usual
covering and automorphism descriptions \cite{GLS1994}, now reduces to a
large-parameter group of Lie type.  We record the numerical exclusions.
A faithful projective degree-$n$ action of $A_5^u$ forces $n\ge2u$.  Lift to the
product of the universal covers.  Each perfect factor acts nontrivially on
some composition factor, since a nonsoluble image cannot act trivially on
all composition factors.  Irreducible modules of a finite direct product
over an algebraically closed field are outer tensors, also in modular
characteristic.  A composition factor on which $a\ge1$ direct factors act
nontrivially has dimension at least $2^a\ge2a$; summing proves the claim.

In alternating degree $d$ there are $\lfloor d/5\rfloor$ disjoint copies
of $A_5$.  In high classical rank $e$ there are at least
$\lfloor(e-1)/5\rfloor$ such factors on ordinary or paired hyperbolic
coordinates; nonsplit forms still have enough coordinate pairs, and
unitary groups can use orthonormal coordinates.  These perfect permutation
groups lie in the projective special group.  Thus the relevant Lie rank
is at most $100(n+1)$, and the alternating degree is bounded as well.
For Lie parameter at most $J^{100}(x)$ in this rank range, or for these
bounded alternating groups, sporadics, and small exceptions, the ordinary
order and automorphism bounds give
$|\operatorname{Aut}(S)|<J^{110}(x)$.  A finite set of this size itself
has format below the avoidance ceiling, contradicting
\eqref{q:almost-simple}.  Hence the parameter satisfies
\begin{equation}\label{q:large-parameter}
                         P>J^{100}(x).
\end{equation}

Suppose its defining prime $p_0$ differs from $p$.  At this parameter the
standard multiplier theorem identifies the simply connected fixed group
$\mathbf H^{F_0}$ of a pinned Steinberg map with the universal central
cover of $S$.  Our projective representation lifts linearly to it and is
projectively nontrivial on noncentral unipotents.  In ordinary
field/diagram type, put $q=P$.  The highest long root $\theta$ is
diagram-stable, also in the simply laced case.  The fixed root elements
have parameters $b=cb^q$ for a fixed $c\ne0$, a one-dimensional additive
$\mathbf F_q$-space.  The fixed coroot elements $\theta^\vee(a)$,
$a\in\mathbf F_q^\times$, act by multiplication by $a^2$.

For a special twisted type write
\[
 F_0=\sigma\operatorname{Fr}_{p_0}^{m},\qquad
 \sigma^2=\operatorname{Fr}_{p_0},\qquad
 q=p_0^{2m+1}=P^2,\qquad p_0\in\{2,3\}.
\]
The highest long and short roots $\theta,\theta'$ are paired by the
special root correspondence, which preserves the positive chamber and
exchanges lengths.  The pinned root-power rules and square-Frobenius
identity hold in each of $B_2,F_4$ in characteristic two and $G_2$ in
characteristic three \citep[Sections~2--3 and Definition~5.4]{DeMedtsNaert2019}.  Their root groups commute because no positive
combination involving both is a root.  In pinned coordinates over the
split prime field,
\[
                 u(b)=x_\theta(b)F_0(x_\theta(b)),
                 \qquad b\in\mathbf F_q,
\]
is a fixed additive copy of $\mathbf F_q$.  On these rank-one parameters
the exponent of $F_0$ is $p_0^r$, where $r=m$ or $m+1$, and the root datum
also gives
\[
 F_0(\theta^\vee(a))=(\theta')^\vee(a^{p_0^r}).
\]
For $a\in\mathbf F_q^\times$, the coroot product
$t(a)=\theta^\vee(a)F_0(\theta^\vee(a))$ is fixed, since the two
coroot factors commute and $F_0^2=\operatorname{Fr}_q$.
The same identity proves fixedness of $u(b)$.  On its $\theta$-coordinate,
$t(a)$ acts by the two root--coroot pairings.  Equivariance makes the
paired coordinate the $F_0$-image of the transformed $\theta$-coordinate.
Thus $t(a)$ sends $u(b)$ to $u(cb)$, where
\[
                   c=a^{2+p_0^r\langle\theta,(\theta')^\vee\rangle}.
\]
The exponent is nonzero at our large $m$ and has absolute value at most
$2+3p_0^{m+1}$.  In ordinary type the scalar subgroup has size at least
$(P-1)/2$; in special type its size is at least
\[
 \frac{P^2-1}{2+3\sqrt{p_0}P}>\frac P{20}.
\]
Thus either case gives a noncentral additive $p_0$-group with a scalar
action subgroup of order at least $P/20$.  It acts diagonalizably in the
linear lift, with a nontrivial character.  Finite-field trace
parametrization shows that the orbit of such a character under the scalar
subgroup is free.  Consequently $n\ge P/20$, contrary to $n\le J(x)$
and \eqref{q:large-parameter}.  The defining characteristic is $p$.

\emph{Claim 3. The outer index satisfies $[\Gamma:S]\le J^3(x)$.}
This bound will control the unions of $S$-cosets used in Claim~4.
Lift to the simply connected universal cover of this large-parameter
group, now with pinned Steinberg map $F$.  The lift is unique by
perfection and irreducible by fullness.  Steinberg's Restriction Theorem,
including its twisted versions, realizes it as the restriction of a
rational simple $L(\lambda)$ of $F$-restricted highest weight
\cite{Steinberg1963,Steinberg1968,Jantzen2003}.  Write its ordinary digits
as $\lambda=\sum_h p^h\lambda_h$.  The Tensor Product Theorem shows
that at most $\log_2n$ digits are nonzero.  Each digit coefficient is less
than $n$: on a highest simple-root face a coefficient $c<p$ gives the
simple rank-one module of dimension $c+1$.

The cyclic subgroup of $\operatorname{Out}(S)$ induced by the pinned
$\operatorname{Fr}_p$ has index at most
$100(100(n+1)+10)^2$.  Every class of this cyclic subgroup realized by
$\Gamma$ preserves the lifted module.  Indeed remove its inner
discrepancy in $S$; a projective intertwining becomes a linear
intertwining on the perfect cover.  For
$F=\operatorname{Fr}_p^f\tau$, with $\tau$ a pinned diagram map, such a
shift preserves the occupied digit positions modulo $f$, by uniqueness
of $F$-restricted restriction.  Wrapping a digit past $f$ only applies a
diagram to it.  Fixing one occupied position bounds the preserving shifts
modulo $f$ by $\log_2n$.  At most three representatives per shift are
needed, since the actual order divides $f\operatorname{ord}(\tau)$.

For special types the unequal root lengths require a different slot count.
For $F=\operatorname{Fr}_p^m\sigma$ with
$\sigma^2=\operatorname{Fr}_p$ and $f=2m+1$, write
\begin{equation}\label{q:special-slot-data}
 \sigma^*\omega_i=p^{c_i}\omega_{i'},\qquad
 c_i,c_{i'}\in\{0,1\},\qquad c_i+c_{i'}=1.
\end{equation}
Here the restricted parametrization is Steinberg's special-factor
version \citep[Sections~11--12, especially Theorem~12.2]{Steinberg1963};
the conversion to the following unequal coefficient boxes and the slot
rotation are deductions from that parametrization and the displayed
root-datum pullback.  The independent universal-cover input used to
lift the projective representation has already been imposed above.
For a short node $a$ paired with a long node $a'$, the pullbacks are
$\sigma^*\omega_a=\omega_{a'}$ and
$\sigma^*\omega_{a'}=p\omega_a$.  On the finite fixed group,
$\sigma$ agrees with $\operatorname{Fr}_p^{m+1}$, by applying $\sigma$
to $F(g)=g$.  The special tensor factorization therefore converts
paired coefficients
\[
 0\le A_a<p^{m+1},\qquad 0\le B_{a'}<p^m
 \quad\longmapsto\quad C_a=A_a+p^{m+1}B_{a'}<p^f
\]
in Steinberg's short-supported parametrization.  The ordinary digit
factors occupy the disjoint ranges $0,\ldots,m$ and $m+1,\ldots,2m$.
This is a bijective change of parameters, simultaneously over all
paired nodes; it gives both existence and uniqueness for the unequal
boxes.  Put
\[
 b_i=m+c_{i'},\qquad b_{i'}=m+c_i,\qquad b_i+b_{i'}=f.
\]
The coefficient at $i$ is less than $p^{b_i}$, so the two node boxes
have $b_i$ and $b_{i'}$ slots.  Their concatenation is a circle of
length $f$.  The first exponent beyond the slots at $i$ is the
$F^*$-image of slot zero at $i'$.

Take shifts $s$ modulo $f$, discarding those for which any occupied
ordinary digit position $h$ satisfies $h+s=m\pmod f$.  At most
$\log_2n$ shifts are discarded.  A nonzero digit newly at
$e=h+s\pmod f$ is kept when $e<m$.  When $e>m$, replace it by the
module with highest weight
\begin{equation}\label{q:crossing-digit}
                         p^{e-m-1}\sigma^*\lambda_h.
\end{equation}
The pinned maps commute, and
\[
 \operatorname{Fr}_p^e=
       \operatorname{Fr}_p^{e-m-1}\sigma F.
\]
Since $F$ fixes the finite group pointwise, this replacement preserves
the finite restriction.  Isogeny pullback preserves irreducibility.
On each node pair the crossing exponent becomes
\[
               e-m-1+c_i=e-b_i
\]
at the other node, exactly rotating the concatenated slots by $s$.
Excluding $e=m$ lets the whole digit be kept or crossed simultaneously
on all its nodes.  There are no overlapping slots.  The resulting
rational tensor has irreducible finite restriction, hence is itself
irreducible, and its sum highest weight is $F$-restricted.  Uniqueness
of restricted simples on the simply connected fixed group forces the
rotation to preserve support whenever the shift is realized by
$\Gamma$.  On a nonempty paired-node cycle at most $2\log_2n$ slots
are occupied, so at most this many rotations preserve its support.
Including the discarded shifts and the bounded diagonal/graph factor
gives, in all types, the ample bound
\begin{equation}\label{q:outer-index}
                           [\Gamma:S]\le J^3(x).
\end{equation}

\emph{Claim 4. There is a positive-height standard point map $\Psi$ of $G$
such that $S=(G^\Psi)'$ and $N_G(S)\le G^\Psi$.}
If several ordinary digits occur, $S$ preserves a nontrivial tensor
splitting.  Its stabilizer in $G$ is proper by the fundamental-module
argument.  At most $[\Gamma:S]$ translates cover $\Gamma$, giving a
proper bounded-format closed set, a contradiction.  The union remains
proper because the ambient connected group variety is irreducible.

For a single digit, untwist its field power, which permutes the pinned
finite points.  The same projective $S$ is then the finite-point image of
a rational module with highest-weight coefficients below $n$.  Let $C$
be its adjoint algebraic domain.  The preimage in $C$ of the projective
image of $G$ is closed, since algebraic-group images are closed, has
format at most $J^{10}(x)$, and contains $(C^F)'$.  To obtain the format
bound use the two common-word action graphs with $y=J^2(x)$, which
covers both ranks, module dimensions, and highest-weight coefficients.
Apply \Cref{q:root-formats} with $y=J^{11}(x)$.  The inequality
$P>J^{100}(x)$ makes the derived-fixed-group order exceed the number
of fixed points allowed in a proper subset, so this preimage is all of
$C$.  The projective image of $C$ cannot be strictly smaller than that
of $G$: its pullback in $G$ would be a proper bounded-format closed
subgroup containing $S$, whose at most $[\Gamma:S]$ translates again
cover $\Gamma$.  The word-graph estimates and \eqref{q:outer-index}
keep both this test and the tensor test below $J^{1000}(x)$.

The adjoint domains are therefore identified on points through these
projective isogenies.  By the isogeny classification, transporting $F$
gives a positive-height standard point-map element $\Psi$.  Up to
inner maps, self-identifications use diagrams and positive or negative
split-field or special powers, and transport conjugates within that
point-map group.  For the dual pair $B_e,C_e$ in characteristic two,
$e\ge3$, the pinned cross isogenies commute with split powers and need
no diagram twist.  We may compose the two projective isogenies with the
adjoint quotient of their common image, which changes nothing on points,
as the maps from the adjoint domains are already point-injective.  Even
if both maps pass through a cross-type target, the cross maps and their
pointwise inverses commute with split powers.  Hence $S=(G^\Psi)'$.
If $g\in G$ normalizes $S$, then $g^{-1}\Psi(g)$ centralizes $S$,
whose inner centralizer is trivial; thus $\Psi(g)=g$.  Since $\Gamma$
normalizes $S$, this proves \eqref{q:sandwich}.
\end{proof}

\subsection{A simultaneous parameter choice}

\begin{definition}\label{q:ceiling-sequence}
Set
\begin{equation}\label{q:x-sequence}
 x_0=10^9,\qquad x_{j+1}=J^{x_j}(x_j)\quad(0\le j<5),
 \qquad M_*=x_5.
\end{equation}
These are ordinary computable integers.
\end{definition}

\begin{theorem}[Size ceiling]\label{q:size-ceiling}
The negative affine-chart construction has a choice of finite parameters
for which its lattice $L$ is nonrepresentable and
\[
                              |L|<M_*.
\]
In particular a nonrepresentable lattice exists among the tables of
cardinality at most the ordinary computable integer $M_*$.
\end{theorem}

\begin{proof}
We verify the quantitative requirements of the structural pruning,
protected-path, and bounded-root arguments.  The parameters are selected
in the order
\[
 x_0\ ;\quad D,q\ ;\quad \ell,r\ ;\quad
 \text{backward trial increments}\ ;\quad m,n\ ;\quad N,g.
\]
This order matters because only the initial graph-pattern round uses a
slot count depending on $n$.  The later-round slot counts and the chart
losses are fixed before the supply of charts is chosen, as required by
\Cref{c:parameter-order}.

\emph{High-type thresholds and structure counts.}
Use $10^8$ as the high cutoff for alternating degree or classical rank;
the remaining ranks are at most $x_0$.  The high automorphism and covering
assertions have only the usual low exceptions.  In the notation of the
direct-system pruning in \Cref{c:direct-block-groups,c:direct-compatibility},
supported classical groups of natural dimension $h$ have full special
order, also after projectivizing, at least
\[
                 s_0^{c h^2-20h},
\]
whereas the full isometry or linear order is at most
$s_0^{c h^2+20h}$, including in small dimensions.  The direct-split
upper bound
\[
                  h!s_0^{c h^2/2+20h}
\]
is strictly smaller than the supported full special order for $h>10^6$.
The linear comparison in \Cref{c:tensor-order} likewise follows from
\[
          h!|F|^{h^2/4+4}<|F|^{h^2/2-20h}
          \qquad(h>10^6).
\]
Perfection, natural irreducibility, and existence of the supported
classical groups already hold there.  In the initial subspace pruning of
\Cref{c:linear-subspaces,c:isometry-subspaces}, the retained small-space
irreducibility cutoff may be $100$.
Completing two such spaces costs dimension at most $400$.  Summing the
endomorphism dimensions at the successive span-growth steps, including
the two spans in linear type, bounds the independent-list requirement by
fewer than $1000^5$ strict steps.

In the direct-block, field, and tensor rounds, the small-special-group,
rotation, and bounded-field cases can all use $10^6$ as the large cutoff
for the indicated dimension or field size.  In particular plane rotations
then have projective order greater than four when required.  Transfer
avoids the ordinary small-special exceptions once total ground dimension
exceeds $10^6$.  This remains so on the complement of a nuisance-plane
tensor slot, by either a high transferred dimension or a large extension
degree.  The compositum-torus test in \Cref{c:isometry-fields}, in total
dimension $4s>10^6$, uses only
\[
                  \frac{|F|^s+1}{2s}>2+|F|+1.
\]
Here $s>250000$, so the inequality holds already at $|F|=2$ and
continues to hold as $|F|$ increases.
The counts here never bound arbitrary matrices over an unbounded field
merely from bounded dimension.  Instead, the order comparisons bound
refinement fibers, comparable split lengths, or the number of distinct
elementwise-commuting simple tensor slots, of which at most two are then
designated.  When matrices inside a bounded plane algebra are counted,
the field is bounded too.  Exceptional matching counts use the bounded
aligned types of a matched pair; only the plane case needs isometry
choices per edge orbit, and that case has bounded field.

For the top-action orbit count in \Cref{c:orbit-bound}, take the six
successive cutoffs
$x_0,J(x_0),\ldots,J^5(x_0)$.  There are at most four class sizes, so
one of the five gaps has no class size strictly between its endpoints.
Call sizes at most its lower endpoint small.  Small sizes are then at
most $J^4(x_0)$, while each large alternating order exceeds $1000$ times
the product of the small factorials, using $J$ at the actual lower
endpoint.  The joint action intersected with the product of the large
alternatings has index at most $2^4\cdot2^2$ times those factorials:
the factors account for parities and the two possible class swaps.
This is enough for the subdirect projection test.  Equality patterns on
at most four large-index occurrences, up to the diagonal
identifications, together with class choices and the individual small
indices, bound the orbit count by, for example,
\begin{equation}\label{q:top-orbit-bound}
                        2^4\cdot15^4\bigl(J^4(x_0)\bigr)^4.
\end{equation}
For fine-cell pairs multiply by at most $(10^6)^2$, from the aggregate
split bounds.  For the tie-chain test on simple-site pairs, fibers of
size at most $2^2$ suffice.  Point cells in the alternating partition
argument have independent supported alternating groups, giving bounded
orbits in every ordered-cell-pair fiber as well.

Invariant partition and subset counts exponentiate these orbit bounds
once.  A further isometry choice on each bounded-field plane edge orbit
costs base at most $(10^6)^4$ per orbit.  A count inside a single
bounded-field plane matrix algebra may even use the entire power-set
size.  In the shared-component ordered-factor arguments, a matrix
centralizer factor of natural degree at most $10^6$ permits at most
$\log_2(10^6)$ nontrivial elementwise-commuting new tensor slots, over
any field.  A large-space factor permits only possibly the whole factor
as a new one there.  Designating one or two components among the
original at most four and these additions has bounded cost.  Therefore
$J^{50}(x_0)$ exceeds all line-count and rich-top tie-chain bounds,
including the small-field and matching choices.  It also bounds the
retained/missing, support, group-type, large/small, and exceptional-case
subdivisions, and the categories and rounds per ancestor.  Recording
even a thousand such tests or small values up to $x_0$ is within this
bound.  No color records an arbitrary degree or an arbitrary matrix.

\emph{The protected path.}
For a nontrivial projective representation in defining characteristic
and high classical natural dimension $v$, the Weyl-orbit lower bound
$v/4$ from \Cref{w:min-degree} defeats both the smallest tensor slots
and the powers $v^{1/s}$, $s\ge2$.  These comparisons already hold at
dimension $10^6$ in \Cref{w:coarsening-count}, leaving fewer than
$10^6$ primes when that argument stops.  The positive-height fixed-order
comparison used in \Cref{w:root-count} follows from the bounded product
factors in high rank: a height step contributes at least $2^{\dim G}$.

For a change of defining characteristic in \Cref{w:characteristic-switch},
the paired linear subgroup has dimension $h\ge v_j/3$.  Its column-shear
characters give $v_i\ge P_j^{h-1}-1$, and hence, for $v_j\ge10^6$,
\begin{equation}\label{q:characteristic-change}
 \log_2v_i\ge\frac{v_j}{5}\log_2P_j.
\end{equation}
A second change would give $v_l\log_2P_l\le5\log_2v_j$.
Using the uniform order bound $|S_l|\le P_l^{3v_l^2}$, we obtain
\[
 \log_2|S_l|\le
 \frac{75(\log_2v_j)^2}{\log_2P_l}
 \le75(\log_2v_j)^2<\frac{2v_j}{5}\le2\log_2v_i,
\]
contradicting the required representation-degree comparison.  The strict
inequality holds at $v_j=10^6$ and thereafter because
$v_j/(\log_2v_j)^2$ increases in this range.  For
\Cref{w:alternating-switch}, the block count and Bochert bound give index
at least $2^{e/7}$ for a proper transitive subgroup in large degree $e$.
The numerical comparison is
\[
                       e!<2^{(\binom e2-1)/7},
\]
since $e!\le e^e$ and $14\log_2e<e-1-2/e$ for $e\ge10^6$.
The classification of $2$-transitive groups has no nonnatural
alternating exception at these degrees.

The word tests in \Cref{w:a-obstruction,w:signed-obstruction} require
only a high constant rank: there must be enough
prefix positions and neighbors, and in the spin case the exponentially
many sign choices must exceed the rank.  Fewer than $100$ pair colors
at a fixed series and a homogeneous set of size ten suffice.  Thus a
high protected path has at most two high alternating nodes and at most
one characteristic switch on its classical subsequence, giving at most
four consecutive same-characteristic classical runs.  Within each run
pigeonhole among the four series, and then apply pair Ramsey to the
shapes of the weight digits in that series.  The total bounds are far
below $J^2(x_0)$.  We may therefore take
\begin{equation}\label{q:Dq}
                  D=J^{1000}(x_0)<x_1,
                  \qquad q=7^D<x_1.
\end{equation}
The extra stages allow a displaced point start and a stage with a
bounded-type parent.  The field order $q$ exceeds every absolute
line-count requirement, and $q-1$ has the distinct prime divisors $2$
and $3$, as required by the affine-chart construction.

\emph{Bounded-type descents and residual length.}
Use $x_1$ as the rank bound in \Cref{q:genericity}.  If its avoidance
hypothesis fails, there is a proper closed $X$ of format at most
$J^{1000}(x_1)$.  By \Cref{w:point-maps}, every point map used in
\Cref{w:witness} has one of the forms
\[
 \operatorname{Int}(g)\tau\operatorname{Fr}_p^h,
 \qquad
 \operatorname{Int}(g)\sigma^\epsilon\operatorname{Fr}_p^h,
 \qquad h\in\mathbf Z,\quad\epsilon\in\{0,1\},
\]
where $\tau$ is a pinned diagram map.  Indeed
$\sigma^2=\operatorname{Fr}_p$ and
$\sigma^{-1}=\sigma\operatorname{Fr}_p^{-1}$.  In the fixed split
models, positive and negative field powers change only equation
coefficients.  The remaining inner, diagram, and special maps have
bounded-format graphs, independently of $h$; special maps occur only
in characteristics two and three.  Thus images and inverse images
are given by bounded formulas after the coefficient twist.  These
point maps are Zariski homeomorphisms, so their images of $X$ are
closed.  Applying the fiberwise equation bound of
\Cref{q:acf-computable} gives equations of degree less than
$J^{1100}(x_1)$ for all the translates and transforms defining $Z$.

For $u$ ambient coordinates and equation degree at most $d_0$, their
entire equation span has dimension at most $\binom{u+d_0}{u}$.
Selecting a basis therefore replaces the possibly infinite
intersection defining $Z$ by that many equations.  This elementary
cost keeps its format below $J^{1200}(x_1)$.  Its conjugation
normalizer is defined by the bounded-format condition
\[
 \forall z\in G\quad
       \bigl(z\in Z\longleftrightarrow gzg^{-1}\in Z\bigr).
\]
Taking equations for this closed subgroup and its monomial-space
dimension keeps the latter below $J^{1400}(x_1)$.  The independent-list
strict descents securing the sandwich are bounded by this dimension;
one or two more list members suffice.

All large-parameter requirements may be imposed by
\begin{equation}\label{q:bounded-cutoff}
                             P>J^{2000}(x_1).
\end{equation}
This follows from the fixed-point versus degree comparison, including
the zero-height fixed loci.  To compare a smaller positive height whose
parameter is at least $\sqrt P$, the usual product formulas, including
the actual central index, give relative errors at most
\[
                          \frac{100(x_1+10)^4}{\sqrt P}.
\]
The minimal leading full-order growth factor per height has integer
square and exceeds the relevant center bound
$z\le10(x_1+10)^2$.  The displayed error is therefore small enough
for the strict order ratio.  For parameters below $\sqrt P$ the gross
order comparison suffices.

All fixed-group orders below the cutoff in
\eqref{q:bounded-cutoff}, as well as bounded alternating, sporadic, and
exceptional-small socle orders, are below $J^{2010}(x_1)$.  Reserving
residual length $J^{2020}(x_1)$ enforces the cutoff by the available
index over the base, both below current path nodes and below tested
bounded-type flats.  In the final bounded-type label estimate $C2^t$,
the sandwich quotient permits $C=2^z$.  Hence the increment
\begin{equation}\label{q:t}
                              t=J(x_1)
\end{equation}
suffices.

\emph{Flag selection and backward pattern increments.}
Only the two continuation colors in \Cref{w:selection} are needed.  Set
\begin{equation}\label{q:ell-r}
                 \ell=J^{2030}(x_1),\qquad r=3\ell.
\end{equation}
Then $Dr<x_2$ and $t<r$.  The tensor-prime and root counts give the
failure union bound
\begin{equation}\label{q:flag-error}
                     2q^{-\ell}D(3Dr+10^6),
\end{equation}
which is smaller than the two-color matching fraction $1/3$.
Indeed $q^\ell>6D(9D\ell+10^6)$ at the displayed choices.

The total number of successive high-pruning rounds in
\Cref{c:pruning-schedule} at one stage is
less than $x_2/100$: these comprise the simultaneous initial subspace
round and the ordered categories and cases relative to at most $D$
ancestors.  Every subsequent round uses at most $x_2$ slots.  Compute
the required increments backwards, beginning with $s=r$.  Suppose the
desired next increment satisfies $s\le y$, where $y\ge x_2$.
The current path rank is at most $Dr$, and the alphabet of coordinate
functionals has size at most $q^s$.  A failure color records an
ancestor, a bounded category and case, and a pullback subset of at most
$q^{Dr+s}$ points.  These use fewer than $J(y)$ colors.  For example,
the subset contribution alone is bounded by
\[
       2^{q^{Dr+s}}\le2^{y^{2y}}<2^{2^{y!}}<J(y),
\]
with ample room for the bounded ancestor and case factors.  The
alphabet and number of slots also fit below $J(y)$.  The definition
of $J$ thus bounds the Hales--Jewett vertical length by $J^2(y)$,
and $J^3(y)$ accommodates the preceding increment including the
horizontal directions.  Fewer than $x_2/100$ backward rounds require
fewer than $3x_2/100$ iterations, strictly below the $x_2$ iterations
defining $x_3$.  In particular the increment sought in the initial
round is below $x_3$.

\emph{Truncation, initial slots, and chart supply.}
Choose
\begin{equation}\label{q:mnNg}
 \begin{gathered}
       m=2x_3,\qquad n=4x_3,\qquad y=J(x_3),\\
       N=J^4(y),\qquad g=J^4(y).
 \end{gathered}
\end{equation}
Path ranks are below $x_2$ and trial increments below $x_3$, so all
path and trial ranks lie strictly below $m$.  The additional at most
$x_2$ independent-list and pair-span allowance still lies below $m$;
residual heights are greater than $2x_3$.  These exceed the reserved
bounded-type residual length.  Initial subspace pruning may ask for
more labels than the interval height, but its argument never requires
the whole list to have a proper span: only pair spans and the bounded
independent lists must fit the truncation.

After $m,n$ are fixed, the initial graph patterns may use $n+x_2$
slots.  This fits below $y=J(x_3)$; the color and alphabet bounds fit
below $J(y)$ as before.  These bounds also cover the direct trial at
a bounded-type parent, without any earlier anchor.  Thus the value
of $N$ in \eqref{q:mnNg} suffices.  The chart loss in any round is
less than $J(y)$, since it is bounded by the ancestor/category/case
colors and two matching failed charts suffice for contradiction.
Keep one successful flat in each surviving chart for the next round.
There are fewer than $x_2/100$ rounds per node, so the choice of $g$
leaves successful charts after all of them.  The supply is private
to each path node, rather than a common pool for the whole tree.
Both $N$ and $g$ are less than $x_4$.

These choices implement the pattern procedure: in a failed chart use
graphs of horizontal-to-vertical linear maps, or affine maps at the
empty start, and color them by the failure data and pullback flat.
The Hales--Jewett family fixes the pullback.  Varying its slots by
scalar multiples of a horizontal functional nonzero at a point of
that flat gives distinct witness flats indexed by an affine parameter
space.  Pair spans add one dimension and contain the whole parameter
line; bounded independent lists add one dimension per new point.
The chart intersection rules therefore turn these patterns into the
line and descent configurations used in the exclusions.  Witnesses
from two distinct failed charts have label intersection exactly the
base group.  Successive rounds retain earlier exclusions because
the later chosen labels contain the earlier successful anchor.
Consequently at each high node a pruned increment-$r$ flat gives the
next almost simple socle with its ancestor protections, and the flag
selection preserves the required root and tensor data.  A bounded-type
parent instead gives the direct bounded-complexity contradiction.
Thus a putative representation either has that contradiction or
forces a protected high path beyond the chosen $D$.  The resulting
finite lattice is nonrepresentable.

\emph{Counting every lattice element.}
Take one leaf for each required leaf attachment.  A chart has at most
\begin{equation}\label{q:flat-count}
                                F_0=2^{q^N}
\end{equation}
flat subsets.  There are $g$ initial charts, branching at most $gF_0$
at each of $D$ generations, and at most $F_0$ leaf attachments per
trunk chart.  Thus the number of all charts is at most
\begin{equation}\label{q:chart-count}
                  C=(D+1)g(1+gF_0)^D(1+F_0).
\end{equation}
Both orientations together have at most $2CF_0$ interior labels, and
there are two common endpoints.  Put $X=x_4$ and $A=J(X)$.  Since
$q,N<X$, we have
\[
                  q^N\le X^X<2^{X!},\qquad
                  F_0<2^{2^{X!}}<A.
\]
Also $D,g<X<A$, so $1+gF_0<A^2$ and
$(D+1)g(1+F_0)<A^4$.  Hence $C<A^{2X+4}$ and
\begin{equation}\label{q:strict-ceiling}
             |L|\le2CF_0+2<A^{2X+7}<J(A)=J^2(X)<x_5.
\end{equation}
The middle inequality follows from $J(A)>2^{2^{A!}}$ and
$A>X\ge10^9$; the last follows from the $x_4$ iterations defining
$x_5$.  This counts the endpoints as well as every interior label and
proves the claimed ceiling.
\end{proof}
